\documentclass{ximera}

\input{../../../preamble.tex}


\definecolor{fvdnavy}{HTML}{071D43}
\definecolor{fvdcyan}{HTML}{20BFC6}
\definecolor{fvdcyanlight}{HTML}{EFFCFB}
\definecolor{fvdmagenta}{HTML}{F10D69}
\definecolor{fvdmagentalight}{HTML}{FFF1F7}
\definecolor{fvdtext}{HTML}{163044}





%=================================================
% Pedagogische omgevingen
% PDF: tcolorbox
% HTML: learning-box
%=================================================

\newenvironment{voorkennis}
{%
    \ifdefined\HCode
        \HCode{%
<section class="learning-box learning-box--prior">
<div class="learning-box__header">
<span class="learning-box__icon" aria-hidden="true"></span>
<span class="learning-box__title">Voorkennis</span>
</div>
<div class="learning-box__content">
        }%
        \begin{itemize}
    \else
       \begin{tcolorbox}[
    enhanced,
    breakable,
    colback=fvdcyanlight,
    colframe=fvdcyan,
    coltext=fvdtext,
    boxrule=0.5pt,
    leftrule=4pt,
    arc=4mm,
    outer arc=4mm,
    left=7mm,
    right=6mm,
    top=5mm,
    bottom=5mm,
    before skip=6mm,
    after skip=6mm,
    title=Voorkennis,
    coltitle=fvdcyan,
    fonttitle=\bfseries\Large,
    attach title to upper={\par\smallskip},
    boxed title style={
        empty,
        left=0mm,
        right=0mm,
        top=0mm,
        bottom=0mm
    },
    overlay={
        \draw[
            fvdcyan,
            line width=0.5pt
        ]
        ([xshift=42mm,yshift=-13mm]frame.north west)
        --
        ([xshift=-7mm,yshift=-13mm]frame.north east);
    }
]
        \begin{itemize}
    \fi
}
{%
    \end{itemize}
    \ifdefined\HCode
        \HCode{%
</div>
</section>
        }%
    \else
        \end{tcolorbox}
    \fi
}


\newenvironment{leerdoelen}
{%
    \ifdefined\HCode
        \HCode{%
<section class="learning-box learning-box--goals">
<div class="learning-box__header">
<span class="learning-box__icon" aria-hidden="true"></span>
<span class="learning-box__title">Leerdoelen</span>
</div>
<div class="learning-box__content">
        }%
        \begin{itemize}
    \else
\begin{tcolorbox}[
    enhanced,
    breakable,
    colback=fvdmagentalight,
    colframe=fvdmagenta,
    coltext=fvdtext,
    boxrule=0.5pt,
    leftrule=4pt,
    arc=4mm,
    outer arc=4mm,
    left=7mm,
    right=6mm,
    top=5mm,
    bottom=5mm,
    before skip=6mm,
    after skip=6mm,
    title=Leerdoelen,
    coltitle=fvdmagenta,
    fonttitle=\bfseries\Large,
    attach title to upper={\par\smallskip},
    boxed title style={
        empty,
        left=0mm,
        right=0mm,
        top=0mm,
        bottom=0mm
    },
    overlay={
        \draw[
            fvdmagenta,
            line width=0.5pt
        ]
        ([xshift=42mm,yshift=-13mm]frame.north west)
        --
        ([xshift=-7mm,yshift=-13mm]frame.north east);
    }
]
        \begin{itemize}
    \fi
}
{%
    \end{itemize}
    \ifdefined\HCode
        \HCode{%
</div>
</section>
        }%
    \else
        \end{tcolorbox}
    \fi
}



\title{Stelsels herbekeken: rang en oplossingsverzamelingen}
\author{Ingmar Herreman}

\begin{document}

% FVD_CHAPTER_DATA_START
% Automatisch gegenereerd uit index.tex.
% Niet handmatig aanpassen.
\fvdchapterdata{1}{Van algebra naar ruimte}{1}
% FVD_CHAPTER_DATA_END








% ============================================================
% ABSTRACT
% ============================================================

\begin{abstract}
Lineaire stelsels, matrixnotatie en de methode van Gauss--Jordan
heb je in het vijfde jaar al bestudeerd.

In dit hoofdstuk halen we die kennis opnieuw naar boven, maar de nadruk
verschuift. Voor omvangrijke reducties gebruiken we ICT. Belangrijker is
dat je een gereduceerde matrix correct kunt lezen, de rang kunt bepalen
en kunt analyseren of een stelsel geen, precies één of oneindig veel
oplossingen heeft.

Die analyse vormt een belangrijke brug van algebra naar ruimtemeetkunde:
oplossingsverzamelingen zullen straks punten, rechten en vlakken in de
ruimte beschrijven.
\end{abstract}

\maketitle


% ============================================================
% VOORKENNIS EN LEERDOELEN
% ============================================================

\begin{voorkennis}
  \item Je herkent een lineair stelsel.
  \item Je kan een stelsel in matrixnotatie schrijven.
  \item Je kent de elementaire rijoperaties van Gauss--Jordan.
  \item Je weet wat een uitgebreide matrix is.
  \item Je hebt eerder matrices met ICT gereduceerd.
\end{voorkennis}

\begin{leerdoelen}
  \item Je kan een lineair stelsel omzetten naar een uitgebreide matrix.
  \item Je kan een uitgebreide matrix met behulp van ICT reduceren.
  \item Je kan uit een gereduceerde matrix de rang aflezen.
  \item Je kan de rang van de coëfficiëntenmatrix en de uitgebreide matrix vergelijken.
  \item Je kan bepalen of een stelsel geen, precies één of oneindig veel oplossingen heeft.
  \item Je kan vrije onbekenden herkennen en een oplossingsverzameling met parameters schrijven.
  \item Je kan je besluit over de oplosbaarheid wiskundig beargumenteren.
  \item Je kan de eerste link leggen tussen een algebraïsche oplossingsverzameling
        en meetkundige objecten in de ruimte.
\end{leerdoelen}


\begin{opmerking}
Dit hoofdstuk is een \textbf{opfrissing}.

Je hebt deze leerstof al gezien. We bouwen de theorie daarom niet opnieuw
vanaf nul op. In het zesde jaar verschuift de aandacht van lange
handberekeningen naar het \textbf{interpreteren van de reductie}.

Voor omvangrijke berekeningen gebruiken we waar mogelijk ICT.
Jij moet de invoer correct kunnen opstellen, het resultaat kunnen lezen
en er een wiskundig besluit uit kunnen trekken.
\end{opmerking}


% ============================================================
% 1. WAAROM STELSELS OPNIEUW?
% ============================================================

\FVDsection{Waarom beginnen we opnieuw met stelsels?}

Een groot deel van de ruimtemeetkunde die volgt, zal uiteindelijk leiden tot
vragen zoals:

\begin{itemize}
    \item Snijden twee meetkundige objecten elkaar?
    \item Bestaat er precies één gemeenschappelijk punt?
    \item Zijn er oneindig veel gemeenschappelijke punten?
    \item Is een combinatie van voorwaarden onmogelijk?
\end{itemize}

Op het eerste gezicht zijn dat meetkundige vragen.

Toch zullen we ze vaak oplossen met \textbf{algebra}.

Wanneer verschillende voorwaarden tegelijk moeten gelden, ontstaat immers
een \textbf{stelsel van vergelijkingen}. De oplossingsverzameling van dat
stelsel vertelt ons dan wat er meetkundig gebeurt.

Daarom frissen we eerst het algebraïsche gereedschap op dat we doorheen
de hele module nodig zullen hebben.


\begin{voorbeeld}[Een voorproefje]

Beschouw de twee vergelijkingen

\[
\begin{cases}
x+y=4,\\
x-y=2.
\end{cases}
\]

Een koppel \((x,y)\) is een oplossing van het stelsel als het
\textbf{beide vergelijkingen tegelijk} doet kloppen.

De oplossing is

\[
(x,y)=(3,1).
\]

Meetkundig stellen beide vergelijkingen een rechte in het vlak voor.
Het punt \((3,1)\) behoort tot beide rechten en is dus hun snijpunt.

\[
\text{\textbf{gemeenschappelijke oplossing}}
\quad\longleftrightarrow\quad
\text{\textbf{gemeenschappelijk punt}}
\]

In de ruimte zullen we precies hetzelfde principe gebruiken voor
rechten en vlakken.
\end{voorbeeld}



% ============================================================
% 2. VAN STELSEL NAAR MATRIX
% ============================================================

\FVDsection{Van een stelsel naar een matrix}

Beschouw een lineair stelsel met \(m\) vergelijkingen en \(n\) onbekenden:

\[
S=
\begin{cases}
a_{11}x_1+a_{12}x_2+\cdots+a_{1n}x_n=b_1,\\
a_{21}x_1+a_{22}x_2+\cdots+a_{2n}x_n=b_2,\\
\vdots\\
a_{m1}x_1+a_{m2}x_2+\cdots+a_{mn}x_n=b_m.
\end{cases}
\]

We kunnen dit compact schrijven als

\[
A\,X=B,
\]

met

\[
A=
\begin{bmatrix}
a_{11}&a_{12}&\cdots&a_{1n}\\
a_{21}&a_{22}&\cdots&a_{2n}\\
\vdots&\vdots&\ddots&\vdots\\
a_{m1}&a_{m2}&\cdots&a_{mn}
\end{bmatrix},
\qquad
X=
\begin{bmatrix}
x_1\\x_2\\\vdots\\x_n
\end{bmatrix},
\qquad
B=
\begin{bmatrix}
b_1\\b_2\\\vdots\\b_m
\end{bmatrix}.
\]


\begin{definitie}
De matrix \(A\) noemen we de \textbf{coëfficiëntenmatrix}.

De matrix

\[
A_b=
\left[
\begin{array}{cccc|c}
a_{11}&a_{12}&\cdots&a_{1n}&b_1\\
a_{21}&a_{22}&\cdots&a_{2n}&b_2\\
\vdots&\vdots&&\vdots&\vdots\\
a_{m1}&a_{m2}&\cdots&a_{mn}&b_m
\end{array}
\right]
\]

noemen we de \textbf{uitgebreide matrix} van het stelsel.
\end{definitie}


\begin{voorbeeld}

Gegeven

\[
S=
\begin{cases}
2x-y+3z=5,\\
3x+2y+2z=4,\\
5x+3y-z=7.
\end{cases}
\]

Dan is

\[
A=
\begin{bmatrix}
2&-1&3\\
3&2&2\\
5&3&-1
\end{bmatrix},
\qquad
B=
\begin{bmatrix}
5\\4\\7
\end{bmatrix},
\]

en de uitgebreide matrix is

\[
A_b=
\left[
\begin{array}{ccc|c}
2&-1&3&5\\
3&2&2&4\\
5&3&-1&7
\end{array}
\right].
\]

\end{voorbeeld}


\begin{opmerking}[Let op de volgorde]

De kolommen van \(A\) horen bij de onbekenden.

Als we de volgorde

\[
x,\ y,\ z
\]

kiezen, moet die volgorde in \textbf{elke vergelijking} worden gerespecteerd.

Zo moet

\[
x+3z=5
\]

in de matrix worden genoteerd als

\[
\begin{bmatrix}
1&0&3&|&5
\end{bmatrix}.
\]

De ontbrekende \(y\)-term heeft coëfficiënt \(0\).
\end{opmerking}



% ============================================================
% 3. GAUSS-JORDAN: WAT MOET JE NOG KENNEN?
% ============================================================

\FVDsection{Gauss--Jordan: wat moet je nog kennen?}

Je kent de methode van Gauss--Jordan al.

Het basisidee blijft hetzelfde: we vervangen de uitgebreide matrix door
een eenvoudiger matrix die \textbf{dezelfde oplossingsverzameling}
beschrijft.

Daarvoor gebruiken we elementaire rijoperaties:

\[
R_i\leftrightarrow R_j,
\]

\[
R_i\longrightarrow \lambda R_i
\qquad (\lambda\neq 0),
\]

en

\[
R_i\longrightarrow R_i+\lambda R_j.
\]


\begin{opmerking}
Deze rijoperaties veranderen de oplossingsverzameling van het stelsel niet.

Daarom mogen we een ingewikkelde uitgebreide matrix vervangen door een
rij-equivalente, eenvoudiger matrix.
\end{opmerking}


\FVDsubsection{Gereduceerde rij-echelonvorm}

Na volledige Gauss--Jordan-reductie verkrijgen we de
\textbf{gereduceerde rij-echelonvorm} van de matrix.

Je herkent die vorm onder andere aan:

\begin{itemize}
    \item alle nulrijen staan onderaan;
    \item elke niet-nulrij heeft een eerste niet-nulelement dat gelijk is aan \(1\);
    \item zo'n eerste \(1\) noemen we een \textbf{spil} of \textbf{pivot};
    \item in een pivotkolom zijn alle andere elementen \(0\);
    \item de pivots schuiven naar rechts wanneer we naar beneden gaan.
\end{itemize}


Bijvoorbeeld:

\[
\begin{bmatrix}
1&0&2&0&-3\\
0&1&-1&0&4\\
0&0&0&1&2\\
0&0&0&0&0
\end{bmatrix}.
\]

De pivots bevinden zich hier in kolom \(1\), \(2\) en \(4\).



% ============================================================
% 4. ICT
% ============================================================

\FVDsection{Reduceren met ICT}

Vanaf nu hoef je een omvangrijk stelsel niet telkens volledig met de hand
te reduceren.

Gebruik de functie van je rekentoestel of software die een matrix naar
haar \textbf{gereduceerde rij-echelonvorm} brengt
(\emph{rref}, \emph{row reduce}, \emph{reduced row echelon form}, \ldots).

De wiskunde begint echter niet wanneer je op ENTER drukt.
Ze begint bij de correcte invoer en vooral bij de interpretatie van de uitvoer.


\begin{opmerking}[ICT is een hulpmiddel, geen besluit]

Bij een oefening van dit hoofdstuk moet je meestal nog steeds zelf:

\begin{enumerate}
    \item de uitgebreide matrix correct opstellen;
    \item de matrix met ICT reduceren;
    \item de pivots herkennen;
    \item de rang bepalen;
    \item vrije onbekenden herkennen;
    \item de oplosbaarheid bepalen;
    \item de oplossingsverzameling correct formuleren;
    \item je besluit kunnen verklaren.
\end{enumerate}
\end{opmerking}


\begin{voorbeeld}[Van ICT-uitvoer naar wiskunde]

Stel dat ICT voor een stelsel met drie onbekenden de volgende
gereduceerde uitgebreide matrix geeft:

\[
\left[
\begin{array}{ccc|c}
1&0&2&3\\
0&1&-1&4\\
0&0&0&0
\end{array}
\right].
\]

De berekening is nu afgelopen, maar de analyse nog niet.

Er zijn twee pivots in de coëfficiëntenmatrix. Dus

\[
\operatorname{rang}(A)=2.
\]

Ook de uitgebreide matrix heeft twee niet-nulrijen:

\[
\operatorname{rang}(A_b)=2.
\]

Er zijn \(n=3\) onbekenden en slechts \(2\) pivotvariabelen.
Er is dus

\[
3-2=1
\]

vrije onbekende.

Omdat

\[
\operatorname{rang}(A)=\operatorname{rang}(A_b)<3,
\]

heeft het stelsel oneindig veel oplossingen.
\end{voorbeeld}



% ============================================================
% 5. RANG
% ============================================================

\FVDsection{De rang van een matrix}

\begin{definitie}
De \textbf{rang} van een matrix is het aantal pivots in haar
gereduceerde rij-echelonvorm.

Equivalent: de rang is het aantal niet-nulrijen in die gereduceerde vorm.
\end{definitie}


\begin{voorbeeld}

\[
A\sim
\begin{bmatrix}
1&0&2&-1\\
0&1&3&4\\
0&0&0&0
\end{bmatrix}.
\]

Er zijn twee pivots, dus

\[
\operatorname{rang}(A)=2.
\]

\end{voorbeeld}


\begin{opmerking}[Belangrijk]

Voor een stelsel moeten we vaak \textbf{twee rangen} onderscheiden:

\[
\operatorname{rang}(A)
\qquad\text{en}\qquad
\operatorname{rang}(A_b).
\]

Dat verschil is essentieel.

De matrix \(A_b\) bevat één extra kolom ten opzichte van \(A\).
Daarom geldt altijd

\[
\operatorname{rang}(A_b)\geq \operatorname{rang}(A).
\]

De rang van de uitgebreide matrix kan dus nooit kleiner zijn dan
de rang van de coëfficiëntenmatrix.
\end{opmerking}



% ============================================================
% 6. OPLOSBAARHEID
% ============================================================

\FVDsection{Wat vertelt de rang over de oplossingen?}

Beschouw een lineair stelsel met \(n\) onbekenden.

Na reductie zijn er slechts drie fundamenteel verschillende mogelijkheden.


\FVDsubsection{Geen oplossing}

Stel dat

\[
\operatorname{rang}(A)<\operatorname{rang}(A_b).
\]

Dan ontstaat in de gereduceerde uitgebreide matrix een rij van het type

\[
\left[
\begin{array}{ccc|c}
0&0&0&c
\end{array}
\right],
\qquad c\neq 0.
\]

Die rij stelt de vergelijking

\[
0=c
\]

voor.

Dat is onmogelijk.

\begin{eigenschap}
Als

\[
\operatorname{rang}(A)<\operatorname{rang}(A_b),
\]

dan is het stelsel \textbf{strijdig} en heeft het geen oplossing.
\end{eigenschap}


\begin{voorbeeld}

\[
A_b\sim
\left[
\begin{array}{ccc|c}
1&0&2&3\\
0&1&-1&4\\
0&0&0&5
\end{array}
\right].
\]

Voor \(A\) zien we twee pivots:

\[
\operatorname{rang}(A)=2.
\]

Voor \(A_b\) verschijnt een extra pivot in de laatste kolom:

\[
\operatorname{rang}(A_b)=3.
\]

De laatste rij betekent

\[
0=5.
\]

Dus

\[
V=\varnothing.
\]

\end{voorbeeld}



\FVDsubsection{Precies één oplossing}

Stel nu dat

\[
\operatorname{rang}(A)=\operatorname{rang}(A_b)=n.
\]

Elke onbekende hoort dan bij een pivotkolom.
Er zijn geen vrije onbekenden.

\begin{eigenschap}
Als

\[
\operatorname{rang}(A)
=
\operatorname{rang}(A_b)
=
n,
\]

dan heeft het stelsel \textbf{precies één oplossing}.
\end{eigenschap}


\begin{voorbeeld}

Voor een stelsel met drie onbekenden geeft ICT:

\[
A_b\sim
\left[
\begin{array}{ccc|c}
1&0&0&2\\
0&1&0&-1\\
0&0&1&4
\end{array}
\right].
\]

Hier geldt

\[
\operatorname{rang}(A)
=
\operatorname{rang}(A_b)
=
3=n.
\]

Dus het stelsel heeft precies één oplossing:

\[
x=2,\qquad y=-1,\qquad z=4,
\]

of

\[
V=\{(2,-1,4)\}.
\]

\end{voorbeeld}



\FVDsubsection{Oneindig veel oplossingen}

Ten slotte kan

\[
\operatorname{rang}(A)
=
\operatorname{rang}(A_b)
<
n.
\]

Het stelsel is dan wel oplosbaar, maar er zijn minder pivots dan onbekenden.
Minstens één onbekende is vrij.

\begin{eigenschap}
Als

\[
\operatorname{rang}(A)
=
\operatorname{rang}(A_b)
<
n,
\]

dan heeft het stelsel \textbf{oneindig veel oplossingen}.
\end{eigenschap}


Het aantal vrije parameters is

\[
\boxed{
n-\operatorname{rang}(A)
}
\]

zolang het stelsel oplosbaar is.


\begin{voorbeeld}

Beschouw een stelsel met drie onbekenden waarvoor

\[
A_b\sim
\left[
\begin{array}{ccc|c}
1&0&2&3\\
0&1&-1&4\\
0&0&0&0
\end{array}
\right].
\]

We hebben

\[
\operatorname{rang}(A)
=
\operatorname{rang}(A_b)
=
2<3.
\]

Er is dus één vrije parameter.

De derde kolom bevat geen pivot, dus \(z\) is vrij.
Stel

\[
z=t,
\qquad t\in\mathbb{R}.
\]

Uit de eerste rij volgt

\[
x+2z=3
\quad\Longrightarrow\quad
x=3-2t.
\]

Uit de tweede rij volgt

\[
y-z=4
\quad\Longrightarrow\quad
y=4+t.
\]

Dus

\[
V=
\left\{
(3-2t,\;4+t,\;t)
\mid t\in\mathbb{R}
\right\}.
\]

We kunnen dezelfde oplossing ook schrijven als

\[
\begin{pmatrix}
x\\y\\z
\end{pmatrix}
=
\begin{pmatrix}
3\\4\\0
\end{pmatrix}
+
t
\begin{pmatrix}
-2\\1\\1
\end{pmatrix},
\qquad t\in\mathbb{R}.
\]

Bewaar vooral die laatste vorm in je achterhoofd.
Ze zal zeer snel opnieuw verschijnen.
\end{voorbeeld}



% ============================================================
% 7. HET BESLISSCHEMA
% ============================================================

\FVDsection{Een vast beslisschema}

Na de reductie hoef je niet te gokken hoeveel oplossingen een stelsel heeft.

Voor een stelsel met \(n\) onbekenden kijk je eerst naar

\[
r=\operatorname{rang}(A)
\qquad\text{en}\qquad
r_b=\operatorname{rang}(A_b).
\]

Daarna geldt:

\[
\boxed{
\begin{array}{ccl}
r<r_b
&\Longrightarrow&
\text{geen oplossing},
\\[2mm]
r=r_b=n
&\Longrightarrow&
\text{precies één oplossing},
\\[2mm]
r=r_b<n
&\Longrightarrow&
\text{oneindig veel oplossingen}.
\end{array}}
\]


In het laatste geval is het aantal vrije parameters

\[
n-r.
\]


\begin{opmerking}[Een goede analyse bevat meer dan één woord]

Schrijf bij een oefening niet alleen:

\[
\text{``oneindig veel oplossingen''.}
\]

Een volledige redenering kan bijvoorbeeld zijn:

\[
\operatorname{rang}(A)
=
\operatorname{rang}(A_b)
=
2<3.
\]

Dus het stelsel is oplosbaar en heeft één vrije parameter.
Daarom zijn er oneindig veel oplossingen.
\end{opmerking}



% ============================================================
% 8. VRIJE ONBEKENDEN
% ============================================================

\FVDsection{Vrije onbekenden en parameters}

Een onbekende waarvan de kolom in de gereduceerde coëfficiëntenmatrix
geen pivot bevat, noemen we een \textbf{vrije onbekende}.

Die onbekende mag vrij worden gekozen.

We vervangen ze meestal door een parameter zoals

\[
t,\quad s,\quad \lambda,\quad \mu.
\]


\begin{voorbeeld}[Twee vrije parameters]

Een stelsel met vier onbekenden \(x,y,z,u\) wordt gereduceerd tot

\[
\left[
\begin{array}{cccc|c}
1&0&2&-1&3\\
0&1&-1&4&2\\
0&0&0&0&0
\end{array}
\right].
\]

Er zijn twee pivots, dus

\[
\operatorname{rang}(A)=2.
\]

Omdat er vier onbekenden zijn,

\[
n-r=4-2=2.
\]

De kolommen van \(z\) en \(u\) bevatten geen pivot.
We kiezen

\[
z=s,
\qquad
u=t.
\]

Dan

\[
x+2z-u=3
\quad\Longrightarrow\quad
x=3-2s+t,
\]

en

\[
y-z+4u=2
\quad\Longrightarrow\quad
y=2+s-4t.
\]

Dus

\[
V=
\left\{
(3-2s+t,\;2+s-4t,\;s,\;t)
\mid s,t\in\mathbb{R}
\right\}.
\]

\end{voorbeeld}



% ============================================================
% 9. HOMOGENE STELSELS
% ============================================================

\FVDsection{Een bijzonder geval: homogene stelsels}

\begin{definitie}
Een lineair stelsel heet \textbf{homogeen} wanneer alle rechterleden
gelijk zijn aan nul:

\[
A\,X=0.
\]
\end{definitie}


Bijvoorbeeld:

\[
\begin{cases}
x+2y-z=0,\\
2x-y+3z=0.
\end{cases}
\]


Een homogeen stelsel heeft altijd minstens de oplossing

\[
x_1=x_2=\cdots=x_n=0.
\]

Die noemen we de \textbf{nuloplossing} of \textbf{triviale oplossing}.


\begin{opmerking}
Een homogeen stelsel kan dus \textbf{nooit strijdig} zijn.

Voor een homogeen stelsel geldt automatisch

\[
\operatorname{rang}(A)
=
\operatorname{rang}(A_b).
\]

Er blijven slechts twee mogelijkheden over:

\[
\operatorname{rang}(A)=n
\]

geeft alleen de nuloplossing, terwijl

\[
\operatorname{rang}(A)<n
\]

oneindig veel oplossingen geeft.
\end{opmerking}


\begin{voorbeeld}

Stel dat een homogeen stelsel met drie onbekenden na reductie geeft

\[
\left[
\begin{array}{ccc|c}
1&2&-1&0\\
0&0&0&0
\end{array}
\right].
\]

Dan

\[
\operatorname{rang}(A)=1
\]

en zijn er

\[
3-1=2
\]

vrije parameters.

Neem

\[
y=s,\qquad z=t.
\]

Dan

\[
x+2s-t=0
\quad\Longrightarrow\quad
x=-2s+t.
\]

Dus

\[
V=
\left\{
(-2s+t,\;s,\;t)
\mid s,t\in\mathbb{R}
\right\}.
\]

\end{voorbeeld}



% ============================================================
% 10. VAN RANG NAAR MEETKUNDE
% ============================================================

\FVDsection{Van algebra naar ruimte}

Hier verschijnt de reden waarom dit hoofdstuk aan het begin van de
ruimtemeetkunde staat.

Beschouw een oplosbaar lineair stelsel in drie onbekenden

\[
x,\ y,\ z.
\]

De oplossingen zijn punten van de driedimensionale ruimte.

Het aantal vrije parameters vertelt dan hoeveel bewegingsvrijheid
de oplossingsverzameling heeft.


\begin{center}
\[
\begin{array}{c|c|c|c}
\operatorname{rang}(A)
&
\text{vrije parameters}
&
\text{aantal oplossingen}
&
\text{meetkundig}
\\
\hline
3&0&1&\text{één punt}\\
2&1&\infty&\text{een rechte}\\
1&2&\infty&\text{een vlak}\\
0&3&\infty&\mathbb{R}^3
\end{array}
\]
\end{center}


\begin{opmerking}
Deze tabel geldt voor een \textbf{oplosbaar} stelsel in drie onbekenden.

Wanneer

\[
\operatorname{rang}(A)<\operatorname{rang}(A_b),
\]

is er geen oplossingsverzameling en dus ook geen gemeenschappelijk
meetkundig object.
\end{opmerking}


\begin{voorbeeld}[Waarom één vrije parameter een rechte oplevert]

We keerden eerder terug naar

\[
V=
\left\{
(3-2t,\;4+t,\;t)
\mid t\in\mathbb{R}
\right\}.
\]

Schrijf dit als

\[
\begin{pmatrix}
x\\y\\z
\end{pmatrix}
=
\begin{pmatrix}
3\\4\\0
\end{pmatrix}
+
t
\begin{pmatrix}
-2\\1\\1
\end{pmatrix}.
\]

Voor \(t=0\) krijgen we het punt

\[
(3,4,0).
\]

Wanneer \(t\) verandert, bewegen we telkens in dezelfde richting

\[
\begin{pmatrix}
-2\\1\\1
\end{pmatrix}.
\]

De oplossingsverzameling vormt dus een \textbf{rechte in de ruimte}.

Binnenkort zullen we deze schrijfwijze herkennen als een
parametrische/vectoriële vergelijking van een rechte.
\end{voorbeeld}



% ============================================================
% 11. DRIE MATRIXUITVOEREN LEZEN
% ============================================================

\FVDsection{Kan je de uitvoer lezen zonder opnieuw te rekenen?}

Dit is de belangrijkste vaardigheid uit dit hoofdstuk.

Stel dat de oorspronkelijke stelsels telkens drie onbekenden bevatten.


\FVDsubsection{Uitvoer A}

\[
\left[
\begin{array}{ccc|c}
1&0&0&2\\
0&1&0&-1\\
0&0&1&5
\end{array}
\right]
\]

We zien:

\[
\operatorname{rang}(A)=3,
\qquad
\operatorname{rang}(A_b)=3.
\]

Omdat \(n=3\),

\[
\operatorname{rang}(A)=\operatorname{rang}(A_b)=n.
\]

Dus er is precies één oplossing:

\[
V=\{(2,-1,5)\}.
\]


\FVDsubsection{Uitvoer B}

\[
\left[
\begin{array}{ccc|c}
1&0&2&4\\
0&1&-3&1\\
0&0&0&0
\end{array}
\right]
\]

We zien:

\[
\operatorname{rang}(A)=2,
\qquad
\operatorname{rang}(A_b)=2.
\]

Er is één vrije parameter:

\[
3-2=1.
\]

Dus het stelsel heeft oneindig veel oplossingen.


\FVDsubsection{Uitvoer C}

\[
\left[
\begin{array}{ccc|c}
1&0&2&4\\
0&1&-3&1\\
0&0&0&1
\end{array}
\right]
\]

Voor de coëfficiëntenmatrix:

\[
\operatorname{rang}(A)=2.
\]

Voor de uitgebreide matrix:

\[
\operatorname{rang}(A_b)=3.
\]

Dus

\[
\operatorname{rang}(A)<\operatorname{rang}(A_b).
\]

De laatste rij betekent

\[
0=1.
\]

Het stelsel is strijdig:

\[
V=\varnothing.
\]


\begin{opmerking}
Uitvoer B en C lijken bijna identiek.

Toch verandert één enkel getal de conclusie volledig:

\[
[0\ 0\ 0\mid 0]
\]

voegt geen nieuwe voorwaarde toe, terwijl

\[
[0\ 0\ 0\mid 1]
\]

een tegenspraak veroorzaakt.

Dat verschil moet je onmiddellijk kunnen herkennen.
\end{opmerking}



% ============================================================
% 12. VEELGEMAAKTE FOUTEN
% ============================================================

\FVDsection{Veelgemaakte fouten}

\begin{opmerking}[Fout 1: alleen het aantal vergelijkingen tellen]

Drie vergelijkingen en drie onbekenden betekenen \textbf{niet automatisch}
dat er één oplossing is.

Vergelijkingen kunnen afhankelijk zijn of elkaar tegenspreken.

De rang bepaalt hoeveel onafhankelijke voorwaarden werkelijk overblijven.
\end{opmerking}


\begin{opmerking}[Fout 2: rang en aantal onbekenden verwarren]

De rang is het aantal pivots.

Het aantal vrije parameters is

\[
n-\operatorname{rang}(A)
\]

alleen wanneer het stelsel oplosbaar is.
\end{opmerking}


\begin{opmerking}[Fout 3: \(0=0\) en \(0=c\) verwarren]

\[
0=0
\]

is altijd waar en levert geen nieuwe voorwaarde.

Maar

\[
0=c,\qquad c\neq0,
\]

is onmogelijk en maakt het volledige stelsel strijdig.
\end{opmerking}


\begin{opmerking}[Fout 4: alleen naar \(\operatorname{rang}(A)\) kijken]

Voor de oplosbaarheid moet je \(A\) en \(A_b\) vergelijken.

Een stelsel met

\[
\operatorname{rang}(A)=2
\]

kan zowel oneindig veel oplossingen als geen enkele oplossing hebben.

Het verschil wordt zichtbaar in

\[
\operatorname{rang}(A_b).
\]
\end{opmerking}


\begin{opmerking}[Fout 5: ICT-uitvoer als eindantwoord beschouwen]

Een gereduceerde matrix is meestal nog geen volledig antwoord.

Je moet ze nog interpreteren en, indien nodig,
de oplossingsverzameling formuleren.
\end{opmerking}



% ============================================================
% 13. WAT MOET BLIJVEN HANGEN?
% ============================================================

\FVDsection{Wat moet je uit deze opfrissing meenemen?}

Voor een stelsel met \(n\) onbekenden is de kern:

\[
\boxed{
\operatorname{rang}(A)<\operatorname{rang}(A_b)
\Longrightarrow
\text{geen oplossing}
}
\]

\[
\boxed{
\operatorname{rang}(A)
=
\operatorname{rang}(A_b)
=
n
\Longrightarrow
\text{één oplossing}
}
\]

\[
\boxed{
\operatorname{rang}(A)
=
\operatorname{rang}(A_b)
<
n
\Longrightarrow
\text{oneindig veel oplossingen}
}
\]

en in het laatste geval:

\[
\boxed{
\text{aantal vrije parameters}
=
n-\operatorname{rang}(A)
}
\]


Maar belangrijker dan de formules is de betekenis:

\begin{center}
\textbf{De rang vertelt hoeveel onafhankelijke voorwaarden aanwezig zijn.}
\end{center}

en

\begin{center}
\textbf{Vrije parameters vertellen hoeveel bewegingsvrijheid
de oplossingsverzameling nog heeft.}
\end{center}

Dat idee nemen we mee naar de ruimtemeetkunde.

Daar zullen oplossingsverzamelingen niet langer alleen rijen getallen zijn.
Ze worden punten, rechten en vlakken in de ruimte.

\end{document}

